\begin{frame}{실습 3: 바위-가위-보 — 순수 전략이 없는 게임}
  \textbf{바위(0) $>$ 가위(1), 가위(1) $>$ 보(2), 보(2) $>$
  바위(0)}인 제로섬 게임(이기면 $+1$, 지면 $-1$, 비기면
  0):
  \begin{itemize}
    \item \textbf{순수 전략의 내쉬 균형이 \textbf{없다}} ---
      상대의 순수 전략(예: 항상 바위)을 알면 항상 보를
      내면 되므로, ``상대가 바위를 낼 것''이라는 예측은
      상대가 알아채고 전략을 바꿈으로써 \textbf{무너지는
      순환}
    \item 이 게임의 내쉬 균형은 \textbf{혼합
      전략(mixed strategy)}으로만 존재 --- 각 행동을
      \textbf{정확히 1/3씩} 독립적으로 섞는 전략
  \end{itemize}
  \vskip0.5em
  \begin{block}{확인하는 수학}
    상대가 $(1/3, 1/3, 1/3)$이면 내가 어떤 순수 전략을 내도
    기대 보상 $= \tfrac13(+1) + \tfrac13(0) + \tfrac13(-1)
    = 0$로 \textbf{동일} --- 어떤 순수 전략도 이탈 이득을
    주지 않음(제로섬이므로 상대도 0).
  \end{block}
  \vskip0.5em
  \begin{itemize}
    \item Q-learning만으로는 이 균형을 \textbf{``행동''으로
      재현하기 어려움} --- $\varepsilon$-greedy는 결국 한
      행동에 \texttt{argmax}가 기울어져 \textbf{잠기려는
      성향}(실제로 특정 행동에 잠기거나, 잠긴 채로 상대에게
      \textbf{착취}당함)
  \end{itemize}
\end{frame}
